% ================================================================
%  Шаблон отчёта для ДЗ №2 «Триаж»
%  Компилируется в Overleaf через XeLaTeX (Menu -> Compiler -> XeLaTeX)
% ================================================================

\documentclass[11pt,a4paper,twocolumn]{article}

% ---- Кодировка и шрифты для русского ----
\usepackage{fontspec}
\usepackage{polyglossia}
\setmainlanguage{russian}
\setotherlanguage{english}
\setmainfont{CMU Serif}
\setsansfont{CMU Sans Serif}
\setmonofont{CMU Typewriter Text}

% ---- Остальные пакеты ----
\usepackage[margin=2cm]{geometry}
\usepackage{graphicx}
\usepackage{booktabs}
\usepackage{array}
\usepackage{hyperref}
\usepackage{xcolor}
\usepackage{amsmath}
\usepackage{amssymb}
\usepackage{caption}
\usepackage{float}
\usepackage{parskip}

% ---- Настройки двухколоночной вёрстки ----
\setlength{\columnsep}{0.7cm}     % расстояние между колонками

\hypersetup{
    colorlinks=true,
    linkcolor=blue!60!black,
    urlcolor=blue!60!black,
    citecolor=blue!60!black,
}

% ---- Заголовок ----
\title{
    \vspace{-2cm}
    \textbf{ДЗ №2: Триаж} \\
    \large Отчёт: битва моделей
}
\author{
    Фамилия Имя, группа \\
    \texttt{kaggle-username: your\_kaggle\_id}
}
\date{\today}

\begin{document}

% ---- Заголовок и вводный блок на всю ширину ----
\twocolumn[
    \begin{@twocolumnfalse}
        \maketitle

        \section*{Как пользоваться этим шаблоном}

        Заполняйте разделы своим содержанием, оставляя структуру. Курсивные подсказки описывают содержание раздела; перед сдачей их необходимо удалить. Критерии оценивания указаны в начале каждого раздела, их можно удалить/оставить на свое усмотрение.

        Объём отчёта: не более 7 страниц в формате PDF. Оценивается содержание, а не объём.

        Все F1 в отчёте — \textbf{macro}, если явно не сказано иное. «Holdout» — ваши 20\% \texttt{train.csv}, отложенные до любых экспериментов.

        \textbf{Что сдать для получения оценки:} PDF этого отчёта с названием \texttt{name\_surname\_2hw.pdf}; \texttt{name\_surname\_2hw.py} или \texttt{name\_surname\_2hw.ipynb}, воспроизводящий финальный сабмит за один запуск (с фиксированными seed).

        \textbf{Итого за задание: 18 место в лидерборде } .

        \textbf{Штраф: минус 2 балла} за утечку информации из валидации или теста в обучение (например, \texttt{fit} энкодера или скейлера на всём train до разбиения на фолды, подбор гиперпараметров по holdout или public). Проверяется по коду.
        \vspace{0.5cm}
    \end{@twocolumnfalse}
]

% ================================================================
\section{Введение}

\textit{Раздел не оценивается отдельно, но является обязательным.}

\textit{В 1--2 абзацах опишите постановку задачи (целевая переменная, признаки, метрика и обоснование её выбора) и структуру данных (число объектов, признаки, классы). Укажите финальную модель, её holdout F1 и результат на public leaderboard. Раздел выполняет роль аннотации.}

% Ваш текст здесь.

% ================================================================
\section{Данные и валидация}
\label{sec:data}

\textbf{Критерии оценивания (3 балла).}
\begin{itemize}
    \item 1 балл — краткий EDA: баланс классов, типы признаков, корреляции, связь признаков с классом.
    \item 1 балл — 3 явно сформулированные гипотезы: что увидели в EDA, что из этого следует для моделей, как будете проверять.
    \item 1 балл — схема валидации (как отложен holdout, какая CV, почему стратифицированная) и проверка, что CV можно доверять: для каждой модели сравните CV F1 и holdout F1 из таблицы~\ref{tab:models} (насколько они отличаются) и проверьте, что если отсортировать модели по CV и по holdout, порядок получится одинаковый.
\end{itemize}
Объём EDA не влияет на оценку; оценивается обоснованность гипотез.

\subsection{Краткий EDA}
% Таблица долей классов, корреляционная матрица или её фрагмент,
% 1--2 графика «признак × класс». Каждой картинке — одно предложение комментария.
%
% Пример вставки картинки в одну колонку:
% \begin{figure}[H]
%     \centering
%     \includegraphics[width=\columnwidth]{figures/classes.png}
%     \caption{Распределение классов ESI.}
%     \label{fig:classes}
% \end{figure}

\subsection{Гипотезы}
\textit{Для каждой гипотезы укажите наблюдение из EDA, ожидаемое следствие для моделей и способ проверки. Результаты проверки гипотез приводятся в разделе~\ref{sec:final}.}

\begin{enumerate}
    \item % Гипотеза 1
\end{enumerate}

\subsection{Схема валидации}
\textit{Опишите способ выделения holdout, схему cross-validation на оставшихся 80\% данных и место препроцессинга в пайплайне. Обоснуйте надёжность полученной оценки качества.}

% Ваш текст.

% ================================================================
\section{Сравнение моделей и Optuna}
\label{sec:models}

\textbf{Критерии оценивания (4 балла).}
\begin{itemize}
    \item 1 балл — пространства поиска для всех четырёх моделей (таблица~\ref{tab:spaces}) с обоснованием границ (почему такие диапазоны, почему log-шкала, почему этот набор гиперпараметров).
    \item 1 балл — итоговая таблица~\ref{tab:models} заполнена полностью: CV, holdout, public, время обучения, время на одном объекте.
    \item 1 балл — анализ Optuna: график сходимости (лучшее значение к trial $N$) и вывод, хватило ли 50 trials.
    \item 1 балл — объяснение, \textbf{почему} модели расположились именно так, через их inductive bias и свойства данных.
\end{itemize}
Объяснение должно опираться на свойства данных и моделей; констатация результата объяснением не является.

\subsection{Пространства поиска}
\begin{table*}[t]
    \centering
    \caption{Пространства поиска Optuna (50 trials на модель).}
    \label{tab:spaces}
    \begin{tabular}{llp{0.5\textwidth}}
        \toprule
        \textbf{Модель} & \textbf{Гиперпараметр} & \textbf{Диапазон и обоснование} \\
        \midrule
        Logistic Regression & \texttt{C}            & \\
                            & \texttt{class\_weight} & \\
        Decision Tree       & \texttt{max\_depth}    & \\
                            & \ldots                & \\
        Random Forest       & \ldots                & \\
        Gradient Boosting   & \ldots                & \\
        \bottomrule
    \end{tabular}
\end{table*}

\subsection{Результаты}
\begin{table*}[t]
    \centering
    \caption{Итоговое сравнение моделей. F1 — macro. Инференс на одном объекте — медиана по 100 вызовам \texttt{predict} на одной строке.}
    \label{tab:models}
    \small
    \begin{tabular}{lccccc}
        \toprule
        \textbf{Модель} & \textbf{CV F1} & \textbf{Holdout F1} & \textbf{Public F1} & \textbf{Обучение, с} & \textbf{1 объект, мс} \\
        \midrule
        Logistic Regression & 0.XXX & 0.XXX & 0.XXX & & \\
        Decision Tree       & 0.XXX & 0.XXX & 0.XXX & & \\
        Random Forest       & 0.XXX & 0.XXX & 0.XXX & & \\
        Gradient Boosting   & 0.XXX & 0.XXX & 0.XXX & & \\
        \bottomrule
    \end{tabular}
\end{table*}

\textit{Лучшие гиперпараметры каждой модели — списком или в приложении.}

\subsection{Анализ Optuna}
% График сходимости, комментарий: где поиск вышел на плато.

\subsection{Почему модели расположились именно так}
\textit{Свяжите результат с inductive bias моделей и с наблюдениями из EDA: какие зависимости в данных линейная модель описать не может, почему одиночное дерево хуже леса, за счёт чего бустинг обгоняет лес.}

% Ваш текст.

% ================================================================
\section{Bias-variance}
\label{sec:bv}

\textbf{Критерии оценивания (2 балла).}
\begin{itemize}
    \item 1 балл — \textbf{график}. Decision Tree, меняется только \texttt{max\_depth} от 1 до 25, остальные гиперпараметры по умолчанию. Для каждой глубины на 5-fold CV посчитаны два числа: macro F1 на train-части фолдов (train F1) и на валидационных фолдах (validation F1), оба усреднены по фолдам. На одном графике две линии; на нём отмечены три зоны — underfitting, оптимальная глубина, overfitting — и подписано, где какая.
    \item 1 балл — \textbf{объяснение}, в котором есть ответы на три вопроса:
    \begin{enumerate}
        \item Что происходит с train F1 и validation F1 при малой и при большой глубине и почему? Например: почему у глубокого дерева train F1 близок к 1, а validation F1 падает?
        \item Как это связано с bias и variance: какая из двух частей ошибки преобладает в зоне underfitting, а какая в зоне overfitting?
        \item Что изменилось бы, если бы обучающая выборка была в несколько раз больше: как сдвинулись бы кривые и оптимальная глубина? Какую часть ошибки больше данных уменьшает, а какую нет?
    \end{enumerate}
\end{itemize}
На графике должны быть подписаны оси и отмечены все три зоны.

% \begin{figure}[H]
%     \centering
%     \includegraphics[width=\columnwidth]{figures/bias_variance.png}
%     \caption{Train и validation F1 от \texttt{max\_depth}.}
% \end{figure}

% Ваш текст.

% ================================================================
\section{Ансамбли}
\label{sec:ensembles}

\textbf{Критерии оценивания (2 балла).}
\begin{itemize}
    \item 1 балл — \textbf{эксперимент}. Random Forest с лучшими гиперпараметрами из Optuna, меняется только \texttt{n\_estimators}: 1, 10, 50, 100, 200, 400. Для каждого значения измерены три величины: macro F1 (на CV или на holdout — укажите, где), время обучения и время предсказания на одном объекте. То же самое сделано для бустинга (число итераций при фиксированном \texttt{learning\_rate}). Результат — на графиках или в таблице; ось \texttt{n\_estimators} удобно брать в логарифмическом масштабе.
    \item 1 балл — \textbf{вывод}, в котором есть ответы на три вопроса:
    \begin{enumerate}
        \item После какого числа деревьев качество леса почти перестаёт расти? Назовите конкретное значение и критерий, по которому вы так решили (например, «прирост меньше 0.005»).
        \item Почему так происходит? Объясните через то, что делает усреднение деревьев: какую часть ошибки (bias или variance) оно уменьшает, а какую нет.
        \item Чем рост \texttt{n\_estimators} в лесу отличается от роста числа итераций в бустинге? Сравните графики двух моделей: как выглядит качество при 1 дереве, как растёт, может ли ухудшиться при слишком большом числе.
    \end{enumerate}
\end{itemize}
В выводе должно быть указано конкретное значение \texttt{n\_estimators}, начиная с которого прирост качества несущественен.


% Ваш текст и графики.

% ================================================================
\section{Macro vs micro F1}
\label{sec:metrics}

\textbf{Критерии оценивания (2 балла).}
\begin{itemize}
    \item 1 балл — \textbf{сравнение метрик}:
    \begin{enumerate}
        \item определения macro F1 и micro F1 (формулами или словами);
        \item посчитайте для одной модели micro F1 и accuracy и сравните их. Объясните, почему в задаче, где у каждого объекта ровно один правильный класс, они получаются одинаковыми, и что из этого следует для редких классов;
        \item таблица~\ref{tab:metrics}: macro F1, micro F1 и F1 по каждому классу на holdout для всех моделей и для константы (всегда предсказывать самый частый класс);
        \item \textbf{пара моделей, которые macro и micro F1 ранжируют по-разному}: модель A лучше B по micro F1, но хуже по macro F1 (или наоборот). Объясните по столбцам F1 по классам, за счёт каких классов так получилось.
    \end{enumerate}
    \item 1 балл — \textbf{подбор гиперпараметров под разные метрики}. Возьмите Decision Tree (он обучается быстрее всех) и запустите Optuna дважды с одинаковым пространством поиска и 50 trials. Различается только метрика, которую Optuna максимизирует: в первый раз micro F1 (\texttt{scoring="f1\_micro"} в \texttt{cross\_val\_score}), во второй — macro F1 (\texttt{scoring="f1\_macro"}). Заполните таблицу~\ref{tab:optuna_metric}: найденные гиперпараметры и метрики на holdout для обоих запусков. В выводе ответьте: чем отличаются найденные деревья, на каких классах разница сильнее всего и что «поощряет» каждая метрика.
\end{itemize}

\begin{table*}[t]
    \centering
    \caption{Метрики на holdout.}
    \label{tab:metrics}
    \small
    \begin{tabular}{lccccccc}
        \toprule
        \textbf{Модель} & \textbf{macro F1} & \textbf{micro F1} & \textbf{ESI-1} & \textbf{ESI-2} & \textbf{ESI-3} & \textbf{ESI-4} & \textbf{ESI-5} \\
        \midrule
        Константа           & & & & & & & \\
        Logistic Regression & & & & & & & \\
        Decision Tree       & & & & & & & \\
        Random Forest       & & & & & & & \\
        Gradient Boosting   & & & & & & & \\
        \bottomrule
    \end{tabular}
\end{table*}

\begin{table*}[t]
    \centering
    \caption{Decision Tree: Optuna максимизирует micro F1 или macro F1. Метрики — на holdout.}
    \label{tab:optuna_metric}
    \small
    \resizebox{\textwidth}{!}{%
    \begin{tabular}{lp{0.22\textwidth}ccccccc}
        \toprule
        \textbf{Метрика Optuna} & \textbf{Лучшие гиперпараметры} & \textbf{macro F1} & \textbf{micro F1} & \textbf{ESI-1} & \textbf{ESI-2} & \textbf{ESI-3} & \textbf{ESI-4} & \textbf{ESI-5} \\
        \midrule
        micro F1 & & & & & & & & \\
        macro F1 & & & & & & & & \\
        \bottomrule
    \end{tabular}%
    }
\end{table*}

\subsection{Пара моделей с разным ранжированием}
% Ваш текст: какие модели, за счёт каких классов (см. таблицу~\ref{tab:metrics}).

\subsection{Подбор под micro F1 и под macro F1}
% Ваш текст: чем отличаются найденные деревья (см. таблицу~\ref{tab:optuna_metric}).

% ================================================================
\section{SHAP и важность признаков}
\label{sec:shap}

\textbf{Критерии оценивания (2 балла).}
\begin{itemize}
    \item 1 балл — SHAP по классам (одинаково ли важны признаки для ESI-1 и ESI-5), dependence plot \texttt{heart\_rate} для ESI-1 с раскраской по \texttt{sbp} с интерпретацией.
    \item 1 балл — сравнение MDI леса, permutation importance на holdout и mean $|\text{SHAP}|$ (таблица~\ref{tab:importance}): где рейтинги расходятся и почему; отдельно про \texttt{queue\_ticket} и про тройку \texttt{sbp} / \texttt{dbp} / \texttt{mbp}; почему SHAP объясняет модель, а не данные.
\end{itemize}
SHAP-значения в многоклассовой задаче вычисляются отдельно для каждого класса; выводы о важности признаков должны это учитывать.

\subsection{Важность по классам}
% График mean |SHAP| по классам.

\subsection{Dependence plot}
% Интерпретация графика и возможное применение наблюдения.

\subsection{Три способа считать важность}
\begin{table}[H]
    \centering
    \caption{Место признака в рейтинге важности.}
    \label{tab:importance}
    \small
    \begin{tabular}{@{}lccc@{}}
        \toprule
        \textbf{Признак} & \textbf{MDI} & \textbf{Perm.} & \textbf{SHAP} \\
        \midrule
        \texttt{chief\_complaint} & & & \\
        \texttt{queue\_ticket}    & & & \\
        \texttt{sbp}              & & & \\
        \ldots                    & & & \\
        \bottomrule
    \end{tabular}
\end{table}

% Ваш текст.

% ================================================================
\section{Финальная модель и выводы}
\label{sec:final}

\textbf{Критерии оценивания (3 балла).}
\begin{itemize}
    \item 1 балл — обоснование выбора финальной модели: почему она, какие эксперименты к этому привели, почему альтернативы хуже; выбор модели при ограничении 10 ms на пациента, подтверждённый замером.
    \item 1 балл — описание как минимум одного неудачного эксперимента с объяснением причин неудачи.
    \item 1 балл — ответы на вопросы подраздела «Выводы», опирающиеся на результаты проведённых экспериментов.
\end{itemize}

\subsection{Выбор финальной модели}
% Ваш текст.

\subsection{Ограничение 10 ms на пациента}
\textit{Укажите, какие модели удовлетворяют ограничению, и приведите результаты замеров.}

% Ваш текст.

\subsection{Что не сработало}
\textit{Опишите эксперимент, который не дал ожидаемого результата: исходную гипотезу, полученный результат и причину расхождения.}

% Ваш текст.

\subsection{Выводы}
\begin{enumerate}
    \item \textit{Какие гипотезы из раздела~\ref{sec:data} подтвердились, а какие нет?}
    \item \textit{Классы ESI упорядочены: ошибка ESI-1 $\to$ ESI-2 менее критична, чем ESI-1 $\to$ ESI-5, однако macro F1 эти ошибки не различает. Какую метрику вы бы предложили вместо неё и что она показывает на ваших моделях?}
    \item \textit{Какие выводы о выборе и настройке моделей можно сделать по результатам работы?}
\end{enumerate}

% Ваш текст.

% ================================================================
% ================================================================
\end{document}